\documentclass[10pt,a4paper]{amsart}
\usepackage[margin=19mm]{geometry}
\usepackage{amsmath,amssymb,mathtools}
\usepackage[scr=rsfs,cal=euler]{mathalfa}
\usepackage{xfrac}
\usepackage[unicode,hidelinks,bookmarks=false]{hyperref}
\newcommand{\ie}{i.e.\ }
\newcommand{\cf}{cf.\ }
\newcommand{\resp}{respectively\ }
\newcommand{\Subsection}{\subsection}
\providecommand{\nobreakdash}{\nobreak\hbox{-}}
\setlength{\emergencystretch}{2em}
\multlinegap0pt
\usepackage{mathtools}
\usepackage{amssymb}
\usepackage[shortlabels]{enumitem}
\usepackage{xcolor}
\newtheorem{defi}{Definition}
\newtheorem{lem}{Lemma}
\usepackage{cleveref}
\crefname{defi}{Definition}{Definitions}
\crefname{lem}{Lemma}{Lemmas}
\newcommand{\Z}{\mathbb{Z}}
\newcommand{\restr}[2]{#1\rvert_{#2}}
\newcommand{\step}[1]{\textbf{#1}}
\DeclareMathOperator{\supp}{supp}
\title{Groupoid Homology and Classifying-Space Homology Are Not Isomorphic}
\author{Luciano Melodia}
\date{Source: arXiv:2602.13375v3 (CC BY 4.0)}
\begin{document}
\maketitle

\noindent\emph{Source and licence.} Groupoid Homology and Classifying-Space Homology Are Not Isomorphic. Version arXiv:2602.13375v3 (CC BY 4.0), distributed by arXiv under the Creative Commons Attribution 4.0 International licence, http://creativecommons.org/licenses/by/4.0/, as recorded in the arXiv licence metadata for this exact version and shown on the abstract page https://arxiv.org/abs/2602.13375v3. Full licence notice: \texttt{licenses/arxiv-2602.13375-CC-BY-4.0.txt}.\par\smallskip

\noindent\textit{Editorial scope.} Counted unit: the article's lem lem:pushforward (source0.tex lines 425--432). The article's complete original proof of it is delivered with it (source0.tex lines 434--492, one \texttt{proof} environment of 59 lines). No mathematics is written, repaired or re-derived: the statement and the proof are the article's own lines, in the article's own notation, and no retained line is substituted or re-flowed.  Retained uncounted as the source-local context the proof needs: lines 410--423.  Nothing was omitted from any retained span, so the package carries no omission record.  No retained line contains a citation, so the wrapper carries no bibliography and no bibliography entry.  No retained line contains a figure environment, an image-inclusion command or drawing code, so no image is delivered and the package declares no asset dependency.  Editorial formatting only, in addition: an AMS \texttt{amsart} preamble that restates the article's own declarations and macros used by the retained lines, namely the environments \texttt{defi}, \texttt{lem} on separate counters, together with the standard packages those lines need.  The retained environments print in this document as Definition 1, Lemma 1 in order of appearance, whereas the article prints the same items with the numbering of its own full document.  The complete native source of the article as distributed in the arXiv e-print for this exact version is bundled unchanged as \texttt{sources/194631/source0.tex}.
\begin{defi}\label{def:Cc}
Let \(Y\) be locally compact Hausdorff and totally disconnected.
Set
\begin{equation}\label{eq:Cc}
C_c(Y,\Z)\coloneqq\{f\colon Y\to\Z\mid f\ \text{locally constant},\ \supp(f)\ \text{compact}\}.
\end{equation}
Since \(f\) is locally constant and \(\Z\) is discrete, \(\{f\neq0\}=f^{-1}(\Z\setminus\{0\})\) is clopen, so \(\supp(f)=\{f\neq0\}\) is a \emph{compact open} subset of \(Y\).

If \(Z\) is locally compact Hausdorff and totally disconnected, \(p\colon Y\to Z\) is a local homeomorphism and \(f\in C_c(Y,\Z)\), define \(p_*f\colon Z\to\Z\) by
\begin{equation}\label{eq:pushforward}
(p_*f)(z)\coloneqq\sum_{y\in p^{-1}(z)}f(y).
\end{equation}
The sum is finite: \(p^{-1}(z)\) is closed (as \(Z\) is Hausdorff and \(p\) continuous) and discrete (as \(p\) is a local homeomorphism), so \(\supp(f)\cap p^{-1}(z)\) is a compact discrete set, hence finite, and only its points contribute.
\end{defi}

\begin{lem}\label{lem:pushforward}
Let \(Y\), \(Z\) and \(W\) be locally compact Hausdorff and totally disconnected, and let \(p\colon Y\to Z\) be a local homeomorphism. Then:
\begin{enumerate}
\item\label{it:push-cc} \(p_*f\in C_c(Z,\Z)\) for every \(f\in C_c(Y,\Z)\),
\item\label{it:push-funct} if \(q\colon Z\to W\) is a local homeomorphism, then \((q\circ p)_*=q_*\,p_*\),
\item\label{it:push-homeo} if \(p\) is a homeomorphism, then \(p_*=(p^{-1})^*\), the pullback \(f\mapsto f\circ p^{-1}\).
\end{enumerate}
\end{lem}

\begin{proof}
Fix \(f\in C_c(Y,\Z)\) and put \(K\coloneqq\supp(f)\), a compact open set with \(\restr{f}{Y\setminus K}=0\).
\begin{itemize}
\item \step{Local constancy of \(p_*f\).}
Fix \(z\in Z\) and enumerate the (finite) fiber of the support,
\begin{equation}\label{eq:fiber-enum}
K\cap p^{-1}(z)=\{y_1,\dots,y_m\},\qquad m\ge 0.
\end{equation}
Because \(p\) is a local homeomorphism, \(f\) is locally constant, and \(Y\) is Hausdorff, we may choose \emph{pairwise disjoint} open sets \(V_1,\dots,V_m\) with \(y_j\in V_j\) such that for each \(j\) the restriction \(\restr{p}{V_j}\colon V_j\to p(V_j)\) is a homeomorphism onto an open set and \(\restr{f}{V_j}\equiv f(y_j)\).
Indeed, a local-homeomorphism chart at \(y_j\) gives the first condition on a neighborhood, intersecting with the clopen set on which \(f\) is constantly \(f(y_j)\) gives the second, and Hausdorffness lets us shrink the finitely many \(V_j\) to be disjoint.
Set \(K'\coloneqq K\setminus(V_1\cup\dots\cup V_m)\).
This is closed in \(K\), hence compact, and contains no point of \(p^{-1}(z)\), since every point of \(K\) over \(z\) is some \(y_j\in V_j\).
Thus \(z\notin p(K')\), and \(p(K')\) is compact, hence closed in \(Z\).
Consequently
\begin{equation}\label{eq:Wprime}
W'\coloneqq\bigl(Z\setminus p(K')\bigr)\cap\bigcap_{j=1}^m p(V_j)
\end{equation}
is an open neighborhood of \(z\) (each \(p(V_j)\ni z\) is open, and for \(m=0\) the empty intersection is read as \(Z\)).
We claim \(p_*f\equiv(p_*f)(z)\) on \(W'\).
Fix \(z'\in W'\).
For each \(j\), since \(z'\in p(V_j)\) and \(\restr{p}{V_j}\) is bijective onto \(p(V_j)\), there is a unique \(y_j'\in V_j\) with \(p(y_j')=z'\), and \(f(y_j')=f(y_j)\).
The \(y_j'\) are distinct as the \(V_j\) are disjoint.
Every \(y\in p^{-1}(z')\) with \(f(y)\neq0\) lies in \(K\), and since \(z'\notin p(K')\) it cannot lie in \(K'\), so it lies in some \(V_j\) and hence equals \(y_j'\).
Therefore
\begin{equation}\label{eq:pf-locally-constant}
\begin{aligned}
(p_*f)(z')
&=\sum_{y\in p^{-1}(z')}f(y)&&\quad\text{by \eqref{eq:pushforward}}\\
&=\sum_{j=1}^m f(y_j')&&\quad\text{the nonzero terms sit at the \(y_j'\)}\\
&=\sum_{j=1}^m f(y_j)&&\quad\text{\(f(y_j')=f(y_j)\)}\\
&=(p_*f)(z).&&\quad\text{by \eqref{eq:pushforward}}
\end{aligned}
\end{equation}
As \(z\in Z\) was arbitrary, \(p_*f\) is locally constant.
\item \step{Compact support of \(p_*f\).}
If \(z\notin p(K)\), then \(p^{-1}(z)\cap K=\varnothing\) and \((p_*f)(z)=0\), so \(\{p_*f\neq0\}\subseteq p(K)\).
Since \(p(K)\) is compact, hence closed, we get \(\supp(p_*f)\subseteq p(K)\), so \(\supp(p_*f)\) is compact.
With local constancy this gives \(p_*f\in C_c(Z,\Z)\).
\item \step{Functoriality.}
Let \(q\colon Z\to W\) be a local homeomorphism.
Then \(q\circ p\) is a local homeomorphism (shrink a chart of \(p\) at each point until its image lies inside a chart of \(q\)), so \((q\circ p)_*\) is defined.
Fix \(w\in W\).
The fibers \(\{p^{-1}(z)\}_{z\in q^{-1}(w)}\) partition \((q\circ p)^{-1}(w)=p^{-1}\bigl(q^{-1}(w)\bigr)\), and \(K\cap(q\circ p)^{-1}(w)\) is finite by the argument of \cref{def:Cc} applied to \(q\circ p\).
Hence the rearrangement
\begin{equation}\label{eq:pf-functorial}
\begin{aligned}
\bigl((q\circ p)_*f\bigr)(w)
&=\sum_{y\in(q\circ p)^{-1}(w)}f(y)&&\quad\text{by \eqref{eq:pushforward}}\\
&=\sum_{z\in q^{-1}(w)}\ \sum_{y\in p^{-1}(z)}f(y)&&\quad\text{partition of \((q\circ p)^{-1}(w)\)}\\
&=\sum_{z\in q^{-1}(w)}(p_*f)(z)&&\quad\text{by \eqref{eq:pushforward}}\\
&=\bigl(q_*(p_*f)\bigr)(w)&&\quad\text{by \eqref{eq:pushforward}}
\end{aligned}
\end{equation}
is a finite reindexing, valid termwise.
Thus \((q\circ p)_*=q_*\,p_*\).
\item \step{Homeomorphisms.}
If \(p\) is a homeomorphism, \(p^{-1}(z)\) is the singleton \(\{p^{-1}(z)\}\), so \((p_*f)(z)=f\bigl(p^{-1}(z)\bigr)=\bigl((p^{-1})^*f\bigr)(z)\). Thus \(p_*=(p^{-1})^*\).\qedhere
\end{itemize}
\end{proof}

\end{document}
